\recipe{Izzy’s (Fiona’s) Brownies}
\label{recipe:brownies}
\makes{16 squares}
\worktime{\SI{30} \minute}
\totaltime{\SI{55} \minute}
\source{Fiona}


\begin{ingredients}

  \item 185g unsalted butter
  \item 185g best dark chocolate
  \item 85g plain flour
  \item 40g cocoa powder
  \item 50g white chocolate
  \item 50g milk chocolate
  \item 3 large eggs
  \item 275g golden caster sugar
\end{ingredients}


\begin{method}
  	Cut butter and dark chocolate into small cubes and tip into a bowl.

Melt the chocolate and butter in a bain marie.Then Leave the melted mixture
to cool to room temperature.

While you wait for the chocolate to cool, position a shelf in the middle of
your oven and turn the oven on to 180C/160C fan/gas 4.

Line the base of a 20cm square tin.

Tip 85g plain flour and 40g cocoa powder into a sieve held over a medium
bowl.

Chop 50g white chocolate and 50g milk chocolate into chunks on a board.

Add 3 eggs and 275g golden caster sugar. Then whisk the eggs and sugar
until they look thick and creamy.

Pour the cooled chocolate mixture over the eggy mousse, then gently fold
together with a rubber spatula.

Resift the flour and cocoa powder into the mixture. Gently fold in this
powder. Stop just before you feel you should, as you don’t want to overdo
this mixing.

Finally, stir in the white and milk chocolate chunks until they’re dotted
throughout.

Pour the mixture into the prepared tin, Put in the oven and set your timer for
25 mins. 

Let it cool down (or not) and eat.

  
\end{method}


\begin{figure}[h]
  	\centering
  	
  	\begin{tikzpicture}
  		% Save the image in a box to get its dimensions
  		\sbox0{\includegraphics[width = 0.6\linewidth]{images/MissingImage.png}}%
  
  		% Draw a clipping rectangle with rounded corners
  		\path[clip, rounded corners=0.5cm] (0,0) rectangle (\wd0,\ht0);
  
  		% Place the image inside the clipped area
  		\node[inner sep=0pt, outer sep=0pt, anchor=south west] at (0,0) {\usebox0};
	\end{tikzpicture}

  \caption*{caption for image with date and information, 1st January 2026}
\end{figure}

\begin{notes}
They’ll keep in an airtight container
for a good two weeks and in the
freezer for up to a month.
\end{notes}

